
\section{补充实验方法}
\label{sec:supp-methods}

\noindent\textbf{构建谱系与部署。}
当前的来源回款代码提交为 \texttt{d99ec92}。归档记录了每个二进制文件的 SHA-256 指纹以及每次运行的配置。当前的九服务续花、回款、存储和短时吞吐系列使用这一来源回款规则。最初的十四服务配对运行使用相同的单义务回款规则，那时批结果标签尚未与直接结果标签分离；当前的跨类型编解码回归覆盖了这一分离。当前存储、内存续花和短时吞吐对照中的成员以 GOGC=200 和 GOMAXPROCS=16 运行；该覆盖设置不应用于委员会、网关或钱包进程。

归档系列保留各自的指纹：运行区间记录基线 \texttt{e3870e1} 与实现 \texttt{a083ff6}；原始续花记录 \texttt{f0c0a274}；资本记录基线 \texttt{8e1fb40} 与实现 \texttt{5807546}；组织路径记录基线 \texttt{bb65af69} 与已验证源码 \texttt{13be856}。归档的批集成使用 \texttt{f82f8cd}；结构性修复计数已在当前代码上重复。来源故障检查点 \texttt{c1ebf2d} 使用已被取代的迟到来源规则，不支持当前的回款论断。

Mac 试验共用一台 M4 Max 主机，16 个 CPU 核和 64\,GiB 内存；结构性微基准运行在 Windows amd64 上。当前的内存链和短时吞吐试验使用内存的服务状态和 NoSync 钱包。回款试验使用磁盘支持的委员会应用状态和 BlockStore，成员与网关在内存中。存储对照仅把成员提交在 bbolt NoSync 与同步事务之间切换。每次运行都从全新创世开始。

另有六次十四服务、磁盘支持的链运行，使用相同的来源回款规则和用户支付的 FUEL。其中三次快速的 100 跳运行耗时 668.960--688.080\,ms；跨跳确认的运行耗时 114.496--118.742\,s。它们属于与正文中九服务结果不同的存储/部署配置。

\noindent\textbf{计时与聚合。}
已认证收款的延迟通常从钱包实际发出第一个 HTTP 请求开始，到收款方验证并在本地原子记录之后结束。资金与来源故障驱动包含少量发送前的构造；它们的比较仅限于同一驱动家族之内。本地 NoSync 事务的完成并不是同步的持久性回执。公共完成与成员完成的时间戳记录的是对已认证执行的观察，而不是委员会的内部提交。发送方的单一单调时钟测量跨站点 ACK 延迟，包含返回路径。收款方可以在前一个发送方收到该 ACK 之前启动其后继付款。提供目标、实际发送速率和含排空的完成速率分别保留。图注区分交易分位数与跨运行的统计量。单次运行内的样本变异性不是对跨运行不确定性的估计。

\noindent\textbf{外部出资的授权控制器。}
资本实验每 200\,ms 采样一次成员可用量。令 $\rho$ 为所声明的每秒 CAL 授权需求，计入每个工作单元中的父请求与子请求；令 $u$ 为认证尚未完成的唯一请求数。控制器设定
\begin{align*}
 b &= 200+100u,\\
 H_{\mathrm{low}} &= \rho\,(0.8\,\mathrm{s})+b,\\
 H_{\mathrm{target}} &= 1.25H_{\mathrm{low}}.
\end{align*}
这里的水位线与突发余量以 CAL 计。当所观察到的最小成员可用量 $A_{\min}$ 低于 $H_{\mathrm{low}}$ 时，所提出的授权额度增量为
\[
 \Delta B=100\left\lceil
 \frac{1.5(H_{\mathrm{target}}-A_{\min})}{100}
 \right\rceil,
\]
并受所配置的授权上限和可用外部资金的约束。同一时刻至多有一次增量在途。其公共命令从外部来源扣款，并原子地向备付金与授权额度入账；其后成员加上由此产生的份额之差，而不清除已有的批准。实验预先提供 $\rho$，而不是估计未知的生产需求。控制器的决定是对实际出资转账的请求，而不是创造备付资本的机制。

\noindent\textbf{资本负载核算。}
每个条件有一次 300-s 的运行，每名成员一个 Worker，并有充足的用户 FUEL。从 3,600 CAL 开始，恒定、速率阶跃和父交易延迟三个条件在外部转账 10,800、18,000 和 19,700 CAL 之后，分别结束于 14,400、21,600 和 23,300 CAL。每个条件完成 6,000 个两跳单元，即 12,000 笔付款，花费 1,128,000 用户 FUEL。相应的固定高备付金对照均能完成；固定低备付金的条件关闭了 6,000 个单元中的 46 个，留下 5,826 个未启动，并记录了 64 笔部分批准的付款。这些类别的单位不同，不能相加为一个付款总数。逐次运行表保留了完整的七个条件的比较，但并不意味着每个条件有三次重复。

\noindent\textbf{身份、资源使用与保留的历史。}
多身份配置使用共享客户端；2,048 个钱包身份并不代表 2,048 个独立的操作系统进程。比较独立输入与十跳付款，三次运行的中位数为每笔付款 12.67 与 15.24 核毫秒，以及每笔付款 59.32 与 71.83\,KiB 的 HTTP 载荷。内存报告把节点总 RSS 与保留的历史，和活动的义务、待处理工作与发件箱分开。因此，历史占用的增加不被报告为付款积压的增加。逐次运行的资源汇总及其归档来源位置包含在归档索引中。

\noindent\textbf{针对特定端点的 Lightning 参考。}
归档的部署使用 LND v0.21.3-beta 和 regtest 中的 Bitcoin Core 31.1，采用同步的 LND 持久化及其 50-ms 通道提交设置。三次全新部署中的每一次都有两个节点和一条直接通道。在 20 次预热转账之后，每次运行使用预先准备的发票，串行地交替发出 100 笔 10,000 聪的付款，以接收方 \texttt{SETTLED} 作为完成事件；报告的 P50/P95 为 282.12/320.17\,ms。另一次 Next 运行以每秒 100 笔的速率发送 300 笔付款，持续 3\,s，使用内存的服务状态和 NoSync 钱包存储。收款方已认证收款的 P50/P95 为 0.9625/1.2075\,ms。负载并发度、持久性和完成语义各不相同，因此这些记录是部署参考，而不是归一化的加速比。


\begin{table}[t]
  \caption{追加式赔付记录与经授权的表示更新（C3）。两者通过相同的经济转移关闭账务；差别在于所存储的历史表示。}
  \label{tab:c3-compare}
  \centering
  \small
  \begin{tabular}{@{}p{0.2\columnwidth}p{0.37\columnwidth}p{0.37\columnwidth}@{}}
    \toprule
    & 追加式记录 & 表示更新（C3） \\
    \midrule
    经济关闭
      & 公共命令中的备付金扣款与义务状态；后继付款不变。
      & 相同的转移与原子步骤，并与一个修订任务一起提交。 \\
    历史表示与原始承诺
      & 原始正文不变；被消费输出的资金来源通过赔付记录的索引找到。
      & 存储的正文指明备付金扣款；交易身份、所有者授权和完整区块身份 $(H_b,P_b)$ 保持不变并经验证。 \\
    实际读者
      & 命令历史与记录的任何读者。
      & 修复构造、安装和本地检查读取修订后的正文；钱包、同步和重放读取原始命令。 \\
    存储与重放
      & 原始区块加记录；重放运行原始命令与赔付事件。
      & 原始区块加修订后的正文、修订和任务；重放运行原始命令与修复事件；同块修复共享分片适配，安装推进到最新的已提交修订。 \\
    进展依赖
      & 赔付记录的公共执行。
      & 赔付与其经验证的修订一起提交，并需要三份适配贡献。 \\
    状态
      & 设计基线；未作为一种模式实现，也未测量。
      & 已实现；成本见正文。 \\
    \bottomrule
  \end{tabular}
\end{table}


\subsection{归档的运行区间系列}
归档的运行区间系列与续花结果相互补充。一个活动组织使用赞助的费用、复用的地址和独立的最终输入。五次各 150,000 笔付款的运行全部完成（图~\ref{fig:operating}）。在目标为每秒 2,000 笔时，两次运行达到含排空的完成速率 1,948.66--1,950.45 笔/秒，收款 P50 为 2.521--2.566\,ms。把目标提高到 2,200，完成速率为 2,103.63--2,106.29 笔/秒。在 2,400 时，达到的速率为 2,174.68，P50/P95 升至 7.174/89.38\,ms，且 4,096 项的待决准入上限记录了 1,635 次命中。运行曲线把达到的完成速率与相应的队列压力配成对。

\begin{figure}[!t]
  \centering
  \includegraphics[width=0.62\columnwidth]{operating-point.pdf}
  \caption{运行点随目标提供负载的变化。上：含排空的完成吞吐（笔/秒）。下：收款延迟（ms）。每个点为一次 150,000 笔付款的运行：目标 2,000 两次，2,200 两次，2,400 一次。P50/P95 汇总每次运行内的交易，而不是跨运行的不确定性。所有服务共用一台主机。}
  \label{fig:operating}
\end{figure}
